\section{Threats to validity\label{sec:threats}}

\paragraph{Construct validity.}
Junior, Mid i Senior to klastry nienadzorowane na obserwowalnej aktywności GitHub, a nie zweryfikowane stanowiska zawodowe czy zmierzone umiejętności. Etykieta ta odzwierciedla względną aktywność w analizowanej populacji, a nie etap kariery. Trzy z dwunastu wskaźników klastrowania - rozpiętość aktywności Agentic-PR, różnorodność agentów i zakres repozytoriów Agentic-PR - opisują bezpośrednio aktywność powiązaną z agentem, a nie niezależne doświadczenie sprzed eksperymentu, bez współpracy z narzędziami AI. Profile są częściowo skonstruowane z tego samego rodzaju zachowania, które później mierzą RQ1 i RQ2. Pozostałe czynniki opisują ogólną aktywność na GitHub niezwiązaną z korzystaniem z agenta. Detekcja botów i agentów opiera się na własnym etykietowaniu AIDev oraz prostym filtrowaniu wzorców loginów, co może błędnie zaklasyfikować nietypowo nazwane konta botów lub ludzi. Dodatkowo pierwotnie wykorzystano flagę zebraną z GraphQL z GitHuba, czy użytkownik jest botem czy nie - okazała się ona znacznie nieskuteczna i w większości przypadków zwracała, że użytkownik jest człowiekiem - choć opis profilu świadczył o czymś innym. Dodatkowo wzbogacanie profili danymi zebrano około 9 miesięcy po punkcie odcięcia AIDev - wskaźniki jak liczba obserwujących czy roczne liczby kontrybucji odzwierciedlają więc stan konta na późniejszą datę zbierania danych, a nie w momencie jego Agentic-PR, więc część profilu aktywności może pochodzić z okresu późniejszego niż wkład, który ma charakteryzować.

\paragraph{Internal validity.}
Projekt badania jest obserwacyjny: żadne z trzech pytań badawczych nie pozwala na wniosek przyczynowy dotyczący wpływu doświadczenia lub korzystania z agenta. Złożoność pull requesta, rodzaj zadania i normy zarządzania repozytorium nie są kontrolowane i prawdopodobnie zakłócają wskaźnik akceptacji. Wynik RQ2, że akceptacja maleje wraz z poziomem aktywności, jest na przykład równie zgodny z hipotezą, że bardziej aktywni deweloperzy podejmują trudniejsze zadania, jak i hipotezą, że agenci gorzej im służą. Parametr $k=3$ zachowano ze względu na interpretowalność, mimo że statystyki nie wskazują ścisłego optimum liczbowego (Sekcja 3). Inny dobór parametru $k$ lub inny algorytm klastrowania mógłby dać inny skład grup, a w konsekwencji inne porównanie wyników. 1755 deweloperów, dla których nie pobrano wskaźnikow z GraphQL, oszacowano na podstawie cech o słabej sile wyjaśniającej (obserwowane $R^2$ pomiędzy $-0,27$ a $0,10$, Sekcja 3) - więc ich przypisanie do profili jest stosunkowo niepewne.

Sprawdzenie wrażliwości typu complete-case - klastrowanie wyłącznie 70 434 deweloperów z rzeczywistymi, nieoszacowanymi danymi, zamiast pełnych 72 189 - wykazało początkowo, że wersja procesu z Sekcji 3 oparta na transformacji kwantylowej jest niestabilna przy tej zmianie populacji: przypisanie do klastrów zgadzało się z klastrowaniem pełnej populacji, dla tych samych nakładających się deweloperów, jedynie przy Adjusted Rand Index = 0,46, a kolejność Junior vs Mid w RQ1 i RQ2 odwróciła się. Prześledziliśmy to do etapu transformacji kwantylowej: ponieważ kilka wskaźników jest silnie wyzerowanych, transformacja kwantylowa mapuje wartości według rangi w obrębie aktualnej próby, więc usunięcie nawet niewielkiej części populacji przesuwa te rangi dla wszystkich pozostałych. Powtórzenie tego samego porównania complete-case z transformacją log(1+x) plus standaryzacją dało Adjusted Rand Index 0,98, czyli praktycznie stabilny wynik, dlatego tej metody użyto w preprocesingu.

Ponownie uruchomiliśmy też RQ1 i RQ2, ograniczając się do 70 434 deweloperów w ramach obecnie przyjętego procesu log(1+x) (RQ3 nie zależy od klastrowania i nie jest tym dotknięte). Oba wyniki były zasadniczo niezmienione (np. RQ1: średnie Junior/Mid/Senior 3,59/42,27/7,18 wobec 3,57/42,12/7,94 w analizie na pełnej populacji; RQ2: 90,3\%/\allowbreak91,8\%/\allowbreak83,4\% wobec 90,2\%/\allowbreak91,8\%/\allowbreak83,2\%), z tym samym wzorcem istotności w porównaniach parami w obu przypadkach. Dla przejrzystości: trzeci kandydujący wynik, mierzący zewnętrzną liczbę komentarzy na Agentic-PR w tych samych trzech profilach, był badany w trakcie prac nad artykułem, ale wykluczony z tego badania po tym, jak nie przeszedł tej samej kontroli - jego kierunek się odwrócił, a istotność statystyczna zniknęła całkowicie pod restrykcją complete-case, w dodatku opierał się na małej i nierówno rozłożonej próbie.

Dodatkowo przeprowadziliśmy cztery dalsze kontrole wrażliwości względem procesu opartego na log(1+x).

1) Alternatywne k: dla k = 4 i 5 trzy główne klastry pozostają bliskie swoim rozmiarom z k = 3, przy czym dodatkowe klastry są wydzielane z Senior jako niewielkie odłamki o wyższej aktywności, a nie w wyniku gruntownej reorganizacji.

2) Powtórzone inicjalizacje: ponowne uruchomienie całego procesu (Isolation Forest i k-średnie) z pięcioma różnymi ziarnami losowości dało Adjusted Rand Index pomiędzy 0,98 a 0,99 względem raportowanego klastrowania, co świadczy o bardzo wysokiej stabilności.

3) Zachowane wartości odstające: pominięcie Isolation Forest w całości daje Adjusted Rand Index = 0,64 względem raportowanego klastrowania, napędzane niemal wyłącznie przez profil Senior, który kurczy się z 13 908 do 1586 deweloperów, gdy najbardziej skrajne profile są izolowane do własnego klastra zamiast być wcześniej wykluczane. Junior i Mid zostają stosunkowo stabilne.

4) Wykluczenie wskaźników pochodzących z AIDev (rozpiętość aktywności Agentic-PR, różnorodność agentów i zakres repozytoriów), pozostawiając jedynie dziewięć pozostałych wskaźników: jest to najmniej stabilna z czterech kontroli, Adjusted Rand Index wyniósł jedynie 0,32, a rozmiary klastrów zmieniają się istotnie (Mid rośnie z 12 705 do 28 329). Ponowne uruchomienie RQ1 i RQ2 przy tym 9-wskaźnikowym klastrowaniu: jakościowy wynik RQ2 się utrzymuje - Senior nadal ma najniższy wskaźnik akceptacji (80,7\% wobec 90,4\% Junior i 90,2\% Mid) - ale kontrast Mid vs Senior w RQ1 jest znacznie słabszy (średnie Agentic-PR 17,6 wobec 20,8, bliskie odwróceniu) niż przy pełnym 12-wskaźnikowym klastrowaniu, co jest spójne z zastrzeżeniem dotyczącym trafności konstruktu podniesionym powyżej - część dużej wielkości efektu RQ1 wynika z tego, że samo klastrowanie jest częściowo zbudowane ze wskaźników korzystania z agenta.

\paragraph{External validity.}
Zbiór danych obejmuje Agentic-PR od pięciu konkretnych agentów kodujących w publicznych repozytoriach GitHub do 1 sierpnia 2025 roku, we wczesnym okresie adopcji agentów. Wzorce użycia, normy akceptacji i zachowania opiekunów wobec kodu generowanego przez agenta mogą się zmieniać w miarę dojrzewania narzędzi i norm społeczności, więc raportowanych związków nie należy przekładać na późniejsze okresy ani na prywatne/firmowe repozytoria bez ponownej weryfikacji. RQ3 jest dodatkowo ograniczone do bazy o ponad 500 gwiazdkach, ponieważ próba porównawcza autorstwa ludzi została zebrana przez autorów AIDev jedynie dla repozytoriów powyżej tego progu popularności; wynik ten opisuje konkretnie repozytoria o wysokiej widoczności i dużej liczbie gwiazdek i może nie uogólniać się na długi ogon mniejszych projektów, które stanowią większość populacji o pełnym zakresie.

\paragraph{Conclusion validity.}
Rozkłady zmiennych wynikowych są silnie nienormalne i wyzerowane, co uzasadniało użycie testów nieparametrycznych i poprawki Bonferroniego w całej pracy (Sekcja 3). Przy licznościach próby rzędu dziesiątek tysięcy nawet bardzo małe różnice mogą osiągnąć istotność statystyczną niezależnie od praktycznego znaczenia; dotyczy to najwyraźniej RQ2 ($\varepsilon^2=0,043$), którego wielkość efektu jest mała, mimo że jego wartość $p$ już nie. RQ1 stanowi wyjątek: jego wielkość efektu ($\varepsilon^2=0,32$) jest duża według konwencjonalnych standardów, więc jego istotność odzwierciedla istotny związek, a nie tylko dużą próbę. RQ3 dodatkowo opiera się na wyraźnie nierównych rozmiarach grup (56\,387 Agentic-PR wobec 2\,286 pull requestów autorstwa ludzi), co dodatkowo ogranicza, jak dużą wagę można przypisać jego oszacowaniu punktowemu, zwłaszcza dla mniejszej grupy.
